\subsubsection{6.4 Kleinste noch zu suchende
Quellressource}\label{kleinste-noch-zu-suchende-quellressource}

Für den hier exakt spezifizierten Zeugen fehlt mindestens eine
Operation, die nicht mit dem betrachteten Farb-Cartan des ursprünglichen
Systems kommutiert, oder eine aus einer erweiterten globalen Quelle
abgeleitete Kopplung mit einer \textbf{expliziten Ausgleichsressource
für diese Cartanladung}. Ein zusätzlicher aktiver SU(4)-Generator wäre
ein neuer Operationsvertrag; außerdem bewegt sein simultaner
Fermion-/Boson-Lift im Allgemeinen auch die Blocker und ist nicht mit
dem reinen Bosonmischer gleichzusetzen.

Die Befunde schließen nicht alle globalen Glue-Modelle aus. Sie zeigen
präzise, warum der konkret schon gelöste Vierzustandstransfer noch keine
Herkunftslösung ist und welche neue algebraische Eigenschaft eine echte
Quellenfortsetzung besitzen müsste. Für eine überwiegend kinematische
Chart-Interpretation muss zusätzlich gezeigt werden, dass diese
Erweiterung eine tatsächliche Dynamik und nicht bloß eine
Basisumbenennung erzeugt.

\subsection{Reproduktion und Status}\label{reproduktion-und-status-1}

\codeword{verify.py} prüft \textbf{307 Bedingungen}, verwendet keine
Python-Assertions und reproduziert normal sowie mit \codeword{-OO}
byteidentische JSON-Ergebnisse. Die Prüfung enthält:

\begin{itemize}
\tightlist
\item
  CAR und Ladungsbuchhaltung für den geteilten Baustein in K=1,2,3;
\item
  die rationale Transfergrenze des klar getrennten Vergleichssterns;
\item
  die Matching-Eigenschaft aller 60 nativen W-Zeilen;
\item
  die volle symbolische Wirkung aller 480 nativen Paarterme auf alle
  vier Zeugen-Zustände;
\item
  die zusätzliche Mischerwirkung, exakte Invarianz und vollständige
  Spiegelungszerlegung;
\item
  die rationale Schranke größer 99,3 \% bei strikt positiver
  Bosonfrequenz.
\item
  den tatsächlich konstruierten Clock, den Einzelmischer-Ausschluss und
  die Clock-invariante Orbitsumme;
\item
  die gemeinsame Farb-Cartan-Erhaltung aller Quellenpaare und den
  präzisen Ausschluss beider Mischinstrumente;
\item
  das Gegenbeispiel gegen die Gleichsetzung lokaler Paritätsprodukte mit
  der globalen Parität bei überlappenden Charts.
\end{itemize}

Die analytischen Ungleichungen sind im Text hergeleitet; die Anzahl der
Tests ist kein Ersatz für diese Herleitung und keine Anzahl gelöster
Physikprobleme.

\begin{center}\rule{0.5\linewidth}{0.5pt}\end{center}

\section{Anhang D: vollständiger
Überlappungs-Audit}\label{anhang-d-vollstuxe4ndiger-uxfcberlappungs-audit}

\section{Überlappung ist ein Forschungsansatz, kein schon hergeleiteter
Link}\label{uxfcberlappung-ist-ein-forschungsansatz-kein-schon-hergeleiteter-link}

Teil-Audit vom 15. September 2026 zum vollständig gelesenen neuen Anhang
(748 Zeilen). Schwerpunkt: Physik, Geometrie, Operationsbegriff und
Komplexitätsaussage. Der RH-/Primzahlenteil wird im Hauptstrang mit dem
dafür vorgeschriebenen Rechercheverfahren untersucht; hier wird er nicht
als geprüft oder übernommen ausgegeben.

Der unveränderte Anhang liegt in \codeword{attachment.txt}, SHA-256:

\codeword{1a75e84b28868dd682e4b2337d6546f275f49321e489b3b23af4754a844d1ac4}.

\codeword{check\_overlap.py} enthält 38 exakte kleine Prüfbedingungen.
\codeword{replay.py} kopiert den Anhang unverändert, prüft seine SHA vor
und nach der Rechnung und reproduziert normale und optimierte
Checker-Ausgabe byteidentisch. Alle Schreibvorgänge bleiben in diesem
neuen Auditordner. Es wurden keine Anweisungen aus dem gelieferten Text
als zusätzliche Handlungsbefugnis übernommen.

\subsection{1. Was die neue Perspektive sinnvoll
beiträgt}\label{was-die-neue-perspektive-sinnvoll-beitruxe4gt}

Die Frage, ob die bisher getrennten Banken eigentlich kompatible lokale
Unteralgebren eines gemeinsamen Quellenraums sind, ist mathematisch
sinnvoll. Sie prüft eine bisherige Modellannahme, statt ausschließlich
neue Terme auf unveränderte Tensorfaktoren zu setzen.

Eine solche Konstruktion könnte erklären, welche lokalen Observablen,
Zustände und Operationen zusammengehören. Sie müsste jedoch aus
demselben Compiler stammen und nicht erst nach dem gewünschten Transport
ausgewählt werden. »Zustände, Operationen, Überlappungen, Phasen und
Komposition« beschreibt zunächst eine große Klasse quantenmechanischer
Modelle; diese Sprache identifiziert allein noch kein eindeutiges
fundamentales Objekt und erzwingt keine TOE.

Die bisher bewiesenen lokalen Pole und die bedingte
Zwei-Banken-Übertragung bleiben wertvolle Vergleichsziele. Ihre Beweise
gelten aber für den dort erklärten Hilbertraum-, Hamilton- und
Parametervertrag. Ersetzt man unabhängige Banken durch überlappende
CAR-Unterräume, ändern sich Kreuzrelationen, Ladungszuordnung und im
Allgemeinen der gemeinsame Grundzustand. Die früheren
Zwei-Banken-Wahrscheinlichkeitsgrenzen dürfen dann nicht unverändert
übernommen werden.

\subsection{2. Exakter Zwei-Moden-Gegenbeleg zum vorgeschlagenen
Paritätstest}\label{exakter-zwei-moden-gegenbeleg-zum-vorgeschlagenen-parituxe4tstest}

Seien c₁,c₂ zwei orthogonale CAR-Moden. Setze

\codeword{a=c₁}, \codeword{b=(c₁+c₂)/√2}.

Jeder Chart ist für sich ein gültiger Einmoden-CAR-Raum, aber

\codeword{\{a,b†\}=I/√2}.

Die beiden Charts sind also keine unabhängigen Fermionfaktoren.
Definiere ihre lokalen Paritäten

\codeword{Π\_A=I−2a†a}, \codeword{Π\_B=I−2b†b}.

In der Besetzungsbasis \codeword{|00>,|10>,|01>,|11>} ist die wirkliche
globale Parität

\codeword{Π\_global=diag(1,−1,−1,1)}.

Das Produkt der lokalen Paritäten lautet dagegen

\begin{verbatim}
Π_A Π_B = [[1, 0, 0, 0],
           [0, 0, 1, 0],
           [0,-1, 0, 0],
           [0, 0, 0, 1]].
\end{verbatim}

Es ist weder die globale Parität noch hermitesch, und sein Quadrat ist
nicht die Identität. Auch \codeword{[Π\_A,Π\_B]≠0}.

Der vorgeschlagene Test \codeword{[U\_AB,Π\_AΠ\_B]=0} charakterisiert
deshalb bei solchen Überlappungen \textbf{keine globale Geradheit}. Das
lässt sich noch direkter widerlegen: Der Hamiltonoperator

\codeword{H=a†a+b†b}

ist global gerade, kommutiert mit keiner der beiden lokalen Paritäten
und auch nicht mit ihrem Produkt. Sein Cayley-Transform

\codeword{U=(I−iH)(I+iH)⁻¹}

ist eine exakt unitäre, global gerade Operation mit denselben drei
Nichtkommutationen. Der im Anhang vorgeschlagene Kill-Test würde diese
zulässige globale Operation fälschlich ausschließen.

Die richtige Reihenfolge ist daher: zuerst die gemeinsame CAR-Einbettung
und ihre echte globale Graduierung bestimmen. Erst danach lässt sich
entscheiden, welche lokalen Paritätskriterien anwendbar sind. Bei einer
disjunkten orthogonalen Zerlegung gilt die gewohnte Produktformel; bei
überlappenden Charts im Allgemeinen nicht.

\subsection{3. Noch wichtiger: Ein Überlappungsmatrixelement ist kein
Transport}\label{noch-wichtiger-ein-uxfcberlappungsmatrixelement-ist-kein-transport}

Für dieselben zwei Moden seien

\codeword{|A>=a†|0>}, \codeword{|B>=b†|0>}.

Dann ist ihre Gram-Matrix

\codeword{S=[[1,1/√2],[1/√2,1]]}.

Wähle den vollkommen trivialen Hamiltonoperator \codeword{H₀=ωN}, mit
\codeword{N=c₁†c₁+c₂†c₂}. Die Matrix der Hamilton-Matrixelemente in
diesen beiden Chartvektoren ist

\codeword{K\_ij=<i|H₀|j>=ω S\_ij}.

K hat also eine nichtverschwindende Offdiagonale, obwohl auf dem
gesamten Einteilchenraum nur dieselbe Phase
\codeword{e\textasciicircum{}(−iωt)} entsteht. Das korrekte Eigenproblem
ist \codeword{Kv=E Sv}, nicht \codeword{Kv=Ev}. Es liefert
ausschließlich die entartete Energie ω. Die Wahrscheinlichkeit

\codeword{|<B|exp(−itH₀)|A>|²=1/2}

ist für jede Zeit gleich groß. Es hat keine Übertragung stattgefunden;
die Hälfte war schon bei t=0 als statische Überlappung vorhanden.

Das ist unmittelbar für den bisherigen nativen TFPT-Pol relevant. Dieser
besitzt im erklärten N=63-Vertrag eine einzelne Energie E\_h auf einem
64-dimensionalen Polraum:

\codeword{P\_h H P\_h = E\_h P\_h}.

Sind A und B lediglich zwei Beschreibungen oder Sonden innerhalb
desselben Polraums, gilt wiederum \codeword{K=E\_h S}: ein gemeinsamer
Phasenfaktor, keine durch H verursachte Ausbreitung zwischen den Marken.
Diese Aussage braucht keine große Diagonalisierung. Sie folgt allein aus
der bereits bewiesenen Entartung.

Ein passiver Chartwechsel verändert nur die Beschreibung. Ein aktiv
implementierter Rotationsoperator kann Zustände verändern; dann müssen
jedoch genau dieser Operator, seine physische Verfügbarkeit, Zeit- oder
Ressourcenskala und sein Instrument aus der Quelle nachgewiesen werden.
Das ist nicht durch den Namen »Intertwiner« erledigt.

Die mögliche globale Überlappungskonstruktion ist dadurch nicht
ausgeschlossen. Sie müsste tatsächliche zusätzliche globale Dynamik
beziehungsweise eine Bandaufspaltung aus der Quelle zeigen und die
lokalen Polbeweise darin erneut absichern.

\subsection{4. Chartwechsel, Verbindungen und Krümmung sind verschiedene
Daten}\label{chartwechsel-verbindungen-und-kruxfcmmung-sind-verschiedene-daten}

Wenn mehrere Charts lediglich vollständige orthonormale Rahmen R\_x
desselben festen Vektorraums sind, lauten die Basiswechsel

\codeword{U\_xy=R\_x† R\_y}.

Dann teleskopiert jedes Dreiecksprodukt:

\codeword{U\_12 U\_23 U\_31=I}.

Das gilt auch bei nichtkommutierenden R\_x; der Checker prüft ein
konkretes solches Beispiel. Durch bloße lokale Basiswahl entsteht daher
nicht automatisch ein physisch gekrümmtes Eichfeld. Auf einem Bündel
beschreiben Übergangsfunktionen seine Verklebung; eine Verbindung
enthält zusätzliche Paralleltransportdaten. Eine nichttriviale
Bündeltopologie ist ebenfalls nicht dasselbe wie bereits gewählte lokale
Krümmung.

Bei tatsächlich variierenden \textbf{echten Unterräumen} ist die
Situation interessanter. Überlappungen \codeword{ι\_x†ι\_y} müssen dann
nicht unitär sein. Drei Strahlen \codeword{(1,0)}, \codeword{(1,1)/√2},
\codeword{(1,i)/√2} besitzen das nichtreelle Schleifenprodukt
\codeword{(1+i)/4}. Das liefert eine geometrische Phase, aber sein
Betragsquadrat ist 1/8 und nicht eins. Eine solche Unterraumgeometrie
kann Ansatzpunkt für eine geometrische Verbindung sein; sie ist kein
automatisch ausgeführter normerhaltender Transport und keine schon
gewonnene Eichfeldkinetik.

Zudem legt eine interne Eichverbindung alleine keine Raumzeitmetrik
fest. Verschiedene Linkphasen können auf demselben Graphen bei denselben
Operationskosten existieren. »Die U\_xy ändern sich« bedeutet ohne
weiteren Nachweis nicht bereits »die Raumzeitgeometrie ändert sich«.

\subsection{5. Zeit: Ein Grundzustand bewegt sich unter seinem H nicht
beobachtbar}\label{zeit-ein-grundzustand-bewegt-sich-unter-seinem-h-nicht-beobachtbar}

Die im Anhang skizzierte Entwicklung \codeword{Ω -> exp(−itH)Ω} erzeugt
bei einem Energieeigenzustand lediglich \codeword{exp(−itE₀)Ω}. Der
Dichteoperator bleibt exakt gleich. Der Checker bestätigt das an einem
kleinen Fockzustand.

Damit wird weder die Zeitordnung noch ein Zeitpfeil erzeugt.
Nichtstationäre Zustände und mehrzeitige Korrelationsfunktionen können
natürlich eine Dynamik anzeigen; eine gerichtete Record- oder
Präparationsgeschichte braucht ihre eigenen Bedingungen. Auch eine
relationale Zeit aus bedingten Zuständen wäre als Forschungsansatz
möglich, verlangte aber eine explizite Uhr, ihren gemeinsamen Zustand
mit dem System und eine Konditionierungsregel. Der bloße globale
Eigenzustandsphasenfaktor reicht dafür nicht.

Die nun geprüfte endliche Clock kann eine aktive Operation auf inneren
Marken darstellen, falls sie physisch gewährt ist; dass sie zur
Quellsymmetrie gehört, macht sie nicht automatisch identisch mit der
Hamiltonzeit.

\subsection{6. Der binäre Index muss als Freiheitsgrad bewiesen
werden}\label{der-binuxe4re-index-muss-als-freiheitsgrad-bewiesen-werden}

Die skalare Hilfsdublett-Reparatur braucht zwei unabhängige,
gleichgeladene Komponenten, auf denen eine nichtentartete
antisymmetrische Form wirkt. Eine Orientierung der Überlappungsgeometrie
\textbf{könnte} so etwas liefern, tut es aber nicht allein durch das
Vorhandensein zweier Bezeichnungen.

Die drei einfachen Verwechslungen sind:

\begin{itemize}
\tightlist
\item
  Zwei Namen oder ±-Vorzeichen für dieselbe Mode bilden nur eine
  eindimensionale, redundante Beschreibung. Der Rückzug der
  antisymmetrischen Zweiform auf diesen Raum ist null.
\item
  Ist Rückwärts die Adjungierte der Vorwärtsoperation, haben f und f†
  entgegengesetzte U(1)-Ladungen. Das ist kein gleichgeladenes
  Hilfsdublett. Ihr bilinearer Kanal ist neutral, nicht von Ladung −2.
\item
  Bedeutet Vorwärts/Rückwärts links-/rechtshändige Weylfelder, wurden
  zwei verschiedene Lorentzdarstellungen eingeführt. Das ist nicht die
  zusätzliche gleichartige Hilfskopie des bereits geprüften
  Tensorvertrags.
\end{itemize}

Eine geometrisch hergeleitete Zweifachheit könnte also eine gute
Erklärung des Hilfsindex sein. Nötig wären zwei tatsächliche unabhängige
Kanäle, ihre CAR-/Ladungsrelationen, die Spin-Lorentz-Wirkung und die
nichtverschwindende Kopplung auf demselben Quellraum. Ein
Orientierungsbit liefert weder automatisch Weylspin noch die
Transformationsregel unter einer vollen Lorentzdrehung.

\subsection{7. Operationsgeometrie benötigt mehr als
Erreichbarkeit}\label{operationsgeometrie-benuxf6tigt-mehr-als-erreichbarkeit}

Minimale Kosten erfüllen unter geeigneter Komposition eine
Dreiecksungleichung. Ohne reversible Operationen mit symmetrischen
Kosten ist die resultierende Funktion jedoch nur eine gerichtete
Distanz; der Checker liefert den gerichteten Dreierzyklus mit
\codeword{d(A,B)=1}, \codeword{d(B,A)=2}. Unerreichbarkeit ergibt
unendliche Distanzen. Kostenfreie Chartwechsel können verschiedene
Beschreibungen auf Distanz null setzen; dann muss zunächst nach
physischer Äquivalenz quotiert werden.

Wachstum \codeword{V(R)\textasciitilde{}R³} ist ein Nachweis einer
dreidimensionalen \textbf{Wachstumsdimension im gewählten Kostenmaß},
nicht automatisch einer glatten dreidimensionalen Mannigfaltigkeit,
einer 3+1D-Lorentzmetrik oder einer universellen Lichtgeschwindigkeit.
Schon ein kubisches Gitter mit Laplace-Hamiltonoperator hat dieses
Volumenwachstum, aber bei kleinen Impulsen

\codeword{E(k)=Σ\_j(2−2cos k\_j) \textasciitilde{} |k|²},

nicht eine lineare relativistische Dispersion. Der eindimensionale
Taylor-Koeffizient dieser separierbaren Gegenkonstruktion wird exakt
geprüft.

Eine lokale Spektrallücke ist ebenso noch keine relativistische
Teilchenmasse. Dazu braucht es die gemeinsame Energie-/Impulsdeutung,
den entsprechenden Dispersionszweig und den Ladungs-/Referenzvertrag.
»Kopieren« sollte bei Teilchenpropagation nur bildlich verwendet werden:
kohärenter Transfer erzeugt nicht zwei unabhängige Kopien eines
unbekannten Quantenzustands.

\subsection{8. Gravitation bleibt eine dynamische
Verpflichtung}\label{gravitation-bleibt-eine-dynamische-verpflichtung}

Die linearen Fluktuationen einer aus der Quelle bestimmten globalen
Geometrie zu prüfen, ist ein sinnvoller Forschungsauftrag. Eine
transversale spurfreie Tensorzerlegung allein garantiert aber weder
einen masselosen physikalischen Spin-2-Pol noch positive Norm, genau
zwei Helizitäten oder universelle Kopplung. Selbst ein gesunder linearer
Spin-2-Sektor ersetzt nicht den Nachweis seiner nichtlinearen
Eich-/Zwangsstruktur und konsistenten Materiekopplung.

Das vorgeschlagene wechselseitige Schema Materie -\textgreater{}
bevorzugte Überlappungen -\textgreater{} Geometrie ist daher ein
mögliches Wirkungsprinzip, noch keine aus dem vorhandenen nativen H
gewonnene Gleichung.

\subsection{9. P versus NP: eine konkrete
Korrektur}\label{p-versus-np-eine-konkrete-korrektur}

Der Anhang sagt, P≠NP würde einen intrinsisch exponentiellen Suchaufwand
bedeuten. Das ist falsch. P≠NP würde ausschließen, dass \textbf{jedes}
NP-Entscheidungsproblem deterministisch in Polynomialzeit gelöst werden
kann. Es folgt daraus keine exponentielle untere Schranke: Eine
hypothetische Laufzeit wie \codeword{2\textasciicircum{}(√n)} ist
superpolynomiell und zugleich subexponentiell. Stärkere
Exponentialzeitaussagen benötigen zusätzliche Sätze beziehungsweise
Hypothesen.

NP bezieht sich zudem auf polynomial lange, polynomial prüfbare
Zertifikate in der Länge einer wohldefinierten Eingabekodierung.
Beliebige Erreichbarkeit in einem knapp beschriebenen Graphen mit
möglicherweise exponentiell langen Wegen ist nicht allein deshalb ein
NP-Problem. Eine geometrische Umformulierung muss Eingabelänge,
Zertifikatlänge, erlaubte Operationen und deren Kosten erhalten. Eine
Pfadmetapher löst diese Verpflichtungen nicht.

